\subsubsection{Resumen}
El World Editor es una herramienta de desarrollo destinada a la creación y edición de mundos dentro del juego. 
Está construido como un cliente de Unity que se comunica con el Core Service a través del API Gateway, siguiendo el mismo 
patrón de acceso que el cliente de juego.

El editor permite a los desarrolladores y diseñadores crear nuevos mundos,
modificar la configuración de zonas existentes, gestionar assets y probar cambios en un entorno controlado antes de su despliegue en producción.

\subsubsection{Estructura de proyecto}
El World Editor se organiza en 3 modulos principales:
\begin{itemize}
  \item \textbf{Core}: Contiene la lógica central del editor, incluyendo gestion de repositorios internos del editor, comunicación con el Core Service, definiciones de interfaces.
  \item \textbf{Gameplay}: Contiene la lógica relacionada con la edición de mundos, como la manipulación de objetos, configuración de zonas y reglas de juego.
  \item \textbf{User Interface}: Incluye todos los elementos visuales del editor, como ventanas, paneles, menús y herramientas de edición.
\end{itemize}
Esta estructura modular permite una clara separación de responsabilidades, facilitando el mantenimiento y la escalabilidad del editor a medida que se agregan nuevas funcionalidades. 
Este mismo patrón de organización se replicó en el cliente y server del juego, siguiendo una arquitectura coherente en todo el sistema.

\subsubsection{Estructura de un Mundo}

Un mundo se compone de tres elementos principales.
\begin{itemize}
  \item \textbf{Zones} (Zonas): Subdivisiones de un mundo, cada una alojada en una instancia de servidor independiente. 
  Esta separación permite escalar horizontalmente y distribuir la carga entre múltiples servidores, mientras el estado 
  del jugador se preserva de forma consistente al cambiar de zona.
\item \textbf{Creatables} (Objetos Creables): Objetos instanciados dinámicamente en tiempo de ejecución, como NPCs, 
enemigos o ítems. A diferencia de los Placeables, no pueden ser posicionados manualmente durante la fase de diseño, 
sino que son generados por la lógica del juego.
\item \textbf{Placeables} (Objetos Colocables): Objetos posicionados manualmente en el editor durante la fase de diseño, 
que permanecen estáticos en tiempo de ejecución. Se clasifican en estructuras (\textit{Structures}), generadores 
(\textit{Spawners}) y objetos diversos (\textit{Miscellaneous Objects}).
\end{itemize}

\subsubsection{Zonas}

En los juegos MMORPG es habitual distribuir el mundo entre múltiples 
servidores para poder sostener entornos de gran escala y una alta cantidad de jugadores simultáneos. Para abordar este 
desafío, cada mundo se divide en zonas independientes, cada una ejecutándose en su propia instancia de servidor dentro 
del clúster. El estado del jugador, incluyendo su inventario y progresión, se mantiene consistente a través de todas 
las zonas de un mismo mundo, permitiendo que el jugador transite entre ellas sin perder información.

Los creadores de mundos pueden crear nuevas zonas dentro del editor, configurando parámetros como la textura
del terreno, la textura del cielo (Skybox) y también configurar como un jugador transporta de una zona a otra mediante
Portales (Portals).

\subsubsection{Placeable Objects (Objetos Colocables)}

Los Placeables son objetos posicionados manualmente en el editor durante la fase de diseño y permanecen estáticos en 
tiempo de ejecución. Se clasifican en tres categorías principales:

\begin{itemize}
  \item \textbf{Structures (Estructuras)}: Objetos estáticos que forman parte del entorno del juego, como edificios, 
  árboles o rocas. Al ser colocados, su posición, rotación y escala pueden editarse mediante \textit{Gizmos} o a través 
  de su panel de configuración. Cada estructura permite configurar sus \textit{Colliders} (Colisionadores), que definen 
  las interacciones físicas con el jugador y otros elementos del juego. Adicionalmente, cualquier estructura puede 
  configurarse como \textit{Shop} (Tienda), habilitando la compra de ítems por parte de los jugadores.

  \item \textbf{Spawners (Generadores)}: Objetos que instancian entidades en tiempo de ejecución, como enemigos o NPCs. 
  Pueden configurarse con parámetros como frecuencia de aparición, rango y tipo de entidad a generar. Actualmente se 
  encuentran disponibles tres tipos: \textit{Enemy Spawner} (Generador de Enemigos), \textit{NPC Spawner} (Generador 
  de NPCs) y \textit{Player Spawner} (Generador de Jugadores).

\item \textbf{Miscellaneous Objects (Objetos Diversos)}: Elementos que no se ajustan a las categorías anteriores pero 
  contribuyen a la composición del mundo. Actualmente se encuentran disponibles dos variantes:
  \begin{itemize}
    \item \textbf{Cofres}: Objetos interactuables que los jugadores pueden abrir para obtener recompensas.
    Configurables con distintos tipos de botín, tiempo de reseteo y apariencia en estado abierto y cerrado.
    \item \textbf{Portales}: Objetos que permiten transportar al jugador a otra ubicación, incluso entre zonas distintas 
    del mundo. Al interactuar con un portal, el jugador es teletransportado a la zona de destino configurada,
    conservando su estado y progresión.
  \end{itemize}

\subsubsection{Creatable Objects (Objetos Creables)}

Los Creatables son objetos instanciados dinámicamente en tiempo de ejecución. A diferencia de los Placeables, no se 
posicionan directamente en el mundo, sino que se configuran y gestionan a través de paneles de edición dedicados.
Actualmente se encuentran disponibles los siguientes tipos:

\begin{itemize}
  \item \textbf{Passive NPC} (Personaje No Jugable Pasivo): Personajes no hostiles con los que el jugador puede 
  interactuar mediante diálogos, como aldeanos o comerciantes. El creador puede asignarles diálogos predefinidos 
  para que el jugador obtenga información, inicie \textit{Quests} o simplemente enriquezca la ambientación del mundo.

  \item \textbf{Aggressive NPC} (Personaje No Jugable Agresivo): Personajes hostiles que atacan al jugador al 
  detectarlo dentro de un rango determinado. Configurables con estadísticas de salud, tipos de armas que utilizan 
  y tabla de botín que sueltan al ser derrotados.

  \item \textbf{Consumable Items} (Objetos Consumibles): Ítems que el jugador puede recoger y utilizar durante 
  el juego, como pociones de salud o comida. Se consumen al ser utilizados y pueden otorgar distintos efectos 
  según su configuración.

  \item \textbf{Weapon Items} (Armas): Ítems que el jugador puede recoger y equipar para atacar NPCs agresivos 
  u otros jugadores. Configurables con tipo de arma (cuerpo a cuerpo o a distancia), daño, alcance, velocidad 
  de ataque y textura.

  \item \textbf{Loot Tables} (Tablas de Botín): Configuraciones que definen qué ítems y qué cantidad de oro 
  se obtienen al interactuar con un elemento del juego. Se asignan a enemigos y cofres, donde cada ítem de la 
  tabla tiene una probabilidad de aparición determinada.

  \item \textbf{Dialogs} (Diálogos): Conjuntos de mensajes asignables a NPCs pasivos que habilitan la 
  interacción con el jugador. Permiten estructurar flujos de conversación y vincular \textit{Quests} a mensajes 
  específicos, posibilitando mostrar distintos diálogos según el progreso del jugador.

  \item \textbf{Quests} (Misiones): Misiones asignables a NPCs que el jugador puede completar para obtener 
  recompensas. Los objetivos disponibles actualmente incluyen eliminar una cantidad determinada de enemigos o 
  interactuar con NPCs específicos. Al completarse, el jugador puede recibir experiencia, oro o ítems especiales.

  \item \textbf{Shop} (Tienda): Configuración asignable a estructuras que habilita la compra de ítems al 
  interactuar con ellas. Soporta dos tipos de productos: ítems consumibles y armas, adquiribles con oro 
  obtenido durante el juego, y cosméticos, adquiribles con gemas.

  \item \textbf{Cosmetics} (Cosméticos): Objetos de personalización estética como ropa, accesorios o skins 
  que no afectan la jugabilidad. Se adquieren con gemas, una moneda premium obtenida con dinero real, y 
  pueden configurarse en tiendas con distintos precios y cantidades disponibles.
\end{itemize}

\subsubsection{Gestión de Mundos}

El editor cuenta con un sistema de guardado y apertura de mundos similar al de otros programas de edición. 
Los mundos se persisten localmente en el equipo del usuario, permitiendo retomar la edición en distintas sesiones de trabajo.

\subsubsection{Publicación de Mundos}

Una vez finalizado el diseño, el creador puede publicar su mundo para que pueda ser accedida mediante el cliente de juego.
El proceso requiere haber iniciado sesión y se realiza desde el panel de publicación. Al confirmar la publicación, 
se inicia una secuencia  de carga hacia el Core Service que sigue el siguiente orden:

\begin{enumerate}
  \item \textbf{World Data}: Los datos estructurales del mundo en formato JSON, que incluyen la configuración 
  de zonas, objetos y sus propiedades.
  \item \textbf{Modelos 3D}: Los modelos utilizados en el mundo, excluyendo aquellos que forman parte del 
  conjunto de assets predeterminados del sistema.
  \item \textbf{Assets de Ítems}: Las imágenes e iconos asociados a los ítems creados por el usuario.
  \item \textbf{Cosméticos}: Los assets visuales correspondientes a los cosméticos configurados en el mundo.
  \item \textbf{Assets de Entorno}: Las texturas de suelo (\textit{ground textures}) y cielos 
  (\textit{skybox textures}) asociadas a cada zona del mundo.
\end{enumerate}

\subsubsection{Gestión de Recursos}

La gestión eficiente de recursos en términos de almacenamiento y rendimiento se resolvió mediante un sistema de 
gestión de assets compuesto por múltiples repositorios internos que abstraen la lógica de persistencia y carga.

% TODO(reference): add GLTFast reference
La importación de modelos 3D en tiempo de ejecución, funcionalidad no soportada nativamente por Unity, se implementó 
mediante la integración de la librería \textit{GLTFast}, que permite cargar modelos en formato \texttt{.glb} durante 
la ejecución del editor, facilitando la incorporación de assets personalizados por parte de los creadores de mundos.


